\documentclass[11pt,a4paper]{article}

% ---------- packages ----------
\usepackage[utf8]{inputenc}
\usepackage[T1]{fontenc}
\usepackage{lmodern}
\usepackage[margin=1in]{geometry}
\usepackage{amsmath,amssymb}
\usepackage{booktabs}
\usepackage{array}
\usepackage{microtype}
\usepackage{xcolor}
\usepackage{graphicx}
\usepackage[most]{tcolorbox}
\usepackage{enumitem}
\usepackage{titlesec}
\usepackage{abstract}
\usepackage[hidelinks]{hyperref}

% ---------- formatting ----------
\setlength{\parskip}{0.35em}
\setlength{\parindent}{0pt}
\titleformat{\section}{\large\bfseries}{\thesection.}{0.6em}{}
\titleformat{\subsection}{\normalsize\bfseries}{\thesubsection}{0.6em}{}
\renewcommand{\abstractnamefont}{\normalfont\bfseries}
\renewcommand{\abstracttextfont}{\normalfont\small}

\definecolor{verdictbg}{HTML}{F3F4F6}
\definecolor{verdictrule}{HTML}{374151}

\newtcolorbox{verdict}[1]{
  colback=verdictbg, colframe=verdictrule, boxrule=0.8pt,
  arc=1.5pt, left=8pt, right=8pt, top=6pt, bottom=6pt,
  title={\bfseries #1}, fonttitle=\bfseries\small, coltitle=verdictrule,
  breakable
}

\newcommand{\status}[1]{{\small\textsc{[#1]}}}
\newcommand{\VR}{\mathrm{VR}}
\newcommand{\half}{t_{1/2}}

% ---------- title ----------
\title{\bfseries From Falsification to a Validated Apparatus:\\[2pt]
\large Mean Reversion in Commodity Calendar Spreads Inside Structurally Trendy Markets}
\author{Adaptive Mean-Reversion (AMR) Research Programme\\
\small Internal Research Dossier --- Compiled Synthesis}
\date{June 2026}

\begin{document}
\maketitle

\begin{abstract}
\noindent
We report the trajectory of a falsification-first research programme over the interval in which its
flagship regime-transition object was killed and a deployment-relevant mean-reversion edge was, for
the first time, validly confirmed. The programme's original object---a pre-reversion ``stabilization''
morphology in the equilibrium residual---is falsified across twelve instruments spanning five habitats
and three sampling resolutions, and the falsification is extended to a sixty-two file cross-habitat
sweep with zero exceptions; a single mechanism, \emph{selection-on-deviation}, reproduces the entire
result with no transition object required. A provenance audit then disqualifies every pre-existing
spread in the data set, and a controlled study shows that even a strictly causal, pre-registered
\emph{rolling} hedge-ratio estimator \emph{manufactures} the variance-ratio structure it purports to
measure---hedge-ratio update noise accounts for $82$--$97\%$ of the differenced-spread variance.
Restricting construction to definitional unit-ratio calendar spreads removes this artifact and
furnishes a real-data positive control. On a natural-gas front calendar the apparatus confirms genuine
short-horizon stochastic mean reversion ($\VR(20)=0.448$, $p=0.005$ against random-walk, GARCH and
moving-average microstructure nulls), corroborated on an independently constructed Brent calendar.
A rolling-local analysis finds the natural-gas signal statistically persistent across nineteen years
(pooled standardized distance $-0.627$, $t=-4.37$, $p\approx 2\times10^{-4}$) and economically
\emph{uneconomic}: a causal mean-reversion book is break-even before costs and negative after. We
maintain throughout a strict separation of three claims---that the \emph{apparatus works}, that the
\emph{signal is real}, and that the edge \emph{clears costs}---and conclude that the first two are
earned and the third is not. We close with the corrected methodological framework, the open problems,
and three forward hypotheses.
\end{abstract}

\noindent\textbf{Keywords:} mean reversion; variance ratio; commodity calendar spreads; theory of
storage; surrogate-relative inference; look-ahead bias; falsification.

\vspace{0.5em}
\hrule
\vspace{0.5em}

\textbf{Status legend.}
\status{frozen} decided, not reopened without justification;
\status{killed/archived} falsified-in-form;
\status{conditional survival} valid in named regimes with a stated kill trigger;
\status{persistent-but-uneconomic} statistically real, does not clear costs;
\status{open} genuinely unresolved;
\status{deferred} deprioritized with a named reopen trigger.
Confidence labels (trustworthiness of evidence): \textsc{low, low--medium, medium, medium--high, high}.

\section{Introduction}

The programme seeks a narrow but genuine edge: to trade mean reversion inside markets that look
structurally trendy, and to \emph{stop}---treating stopping as a successful outcome---if the edge is
not there. This dossier covers the interval bounded by two events. It opens at the falsification of the
programme's flagship regime-shift object and closes at the present research frontier. The narrative is
a single arc: \emph{diagnosis}, \emph{falsification}, \emph{reconstruction}, \emph{corrected
framework}, and \emph{forward path}.

The organising discipline of the report is a separation that the field routinely blurs. A research
result can be a success in three different and non-interchangeable senses:

\begin{description}[leftmargin=1.4em,itemsep=2pt]
  \item[Research win.] The measurement apparatus works: it has power, selectivity, a calibrated
  false-positive rate, and the ability to discriminate genuine mean reversion from microstructure
  noise, demonstrated on controls of known construction.
  \item[Empirical win.] The signal is real: a specific instrument separates from a matched surrogate
  in the mean-reverting direction, robustly across regimes.
  \item[Deployment win.] The edge clears costs: a causal trading rule has positive net expectancy
  after all-in transaction costs over a forward hold-out.
\end{description}

As of this writing the programme holds a research win and an empirical win. It does \emph{not} hold a
deployment win. Every claim below is filed under exactly one of these headings.

The principal methodological commitments are inherited and unrelaxed: at time $t$ only information
available at $t$ may enter a causal computation; every statistic is read \emph{surrogate-relative}
(against a matched null run through identical extraction), never against a fixed constant or the
instrument's own history; hypotheses, windows, and decision rules are fixed before data are touched;
and a falsified idea is named and closed rather than silently revived.

\section{The Falsification of the Pre-Reversion Morphology}
\label{sec:statet}

\subsection{The object and its pre-registered signature}

The original Layer-2 object posited that an extended deviation from a latent equilibrium price
$\mu^*$ \emph{loses momentum and stabilizes before reverting}, leaving a measurable pre-reversion
morphology. This was operationalized---deliberately as an observational existence question, not a
detector---on the causal one-step innovation stream of an anchored Kalman equilibrium filter. The
signature was fixed in advance: in high-$|z|$ causal pre-windows (selected by a symmetric partition on
the causal standardized deviation, $\theta\in\{1.0,1.5,2.0\}$, window $W=30$), relative to a matched
null, all three of the following would hold:
\begin{enumerate}[label=(D\arabic*),leftmargin=2.4em,itemsep=1pt]
  \item innovation variance $\downarrow$ (residual stabilization);
  \item innovation autocorrelation $\mathrm{ACF}(1)\downarrow$ (deviation losing momentum);
  \item directional efficiency $\downarrow$ (price chops rather than continues).
\end{enumerate}
A ``signature'' required all three negative. The claim was distributional---about an ensemble, never a
single bar---so that its meaning did not depend on \emph{when} it was read.

\subsection{The method was validated before it was trusted}

Synthetic calibration established that the measurement was not blind: it produced no false positive on
ordinary Ornstein--Uhlenbeck (OU) or random-walk mean reversion (all standardized mean differences
$|d|\leq 0.15$), and it fired loudly and in the correct direction on a clean stabilize-then-revert
positive control ($d\approx -1.9$ for innovation variance, $-2.1$ to $-2.5$ for $\mathrm{ACF}(1)$,
$-0.9$ for directional efficiency). A first, mis-built positive control was rejected and recorded
rather than quietly discarded. The method's confidence entered the real-data phase at
\textsc{medium--high}; the existence of the object itself remained \textsc{untested}.

\subsection{The kill}

The decisive instrument was committed in advance: a canonical mean-reverting equity pair spread. The
committed rule stated that directional continuation in that instrument, like the rest of the cohort,
would meet the kill condition. It did. Across twelve instruments spanning five habitats, two to three
resolutions, two comparison methods and three thresholds, directional efficiency was \emph{positive in
every one}---high-deviation windows were \emph{more} directionally efficient, the exact opposite of the
pre-registered stabilization. The within-instrument, confound-controlled comparison reproduced the
sign, ruling out a null-calibration artifact.

A single null reproduces the entire result with no transition object: \emph{selection-on-deviation}.
Conditioning on $|z|\geq\theta$ samples paths that arrived at an extreme \emph{by directional travel};
those pre-windows are directionally efficient by construction. The effect is the Lo--MacKinlay
\cite{LoMacKinlay1990} cross-autocovariance selection trap, and continuation around extremes is the
expected result at these horizons in the momentum literature \cite{GeorgeHwang2004,JT1993}. The
result is textbook, not anomalous.

A subsequent sweep over sixty-two files and twenty-plus instruments---equities, commodities, foreign
exchange and rates at three resolutions---returned sixty-two rejections, zero confirmations, zero
inconclusive and zero noise, with the relevant descriptors universally and substantially wrong-signed.

\begin{verdict}{Verdict --- Section \ref{sec:statet}}
\small
The pre-reversion stabilization morphology is \status{killed/archived} \textbf{(structural;
confidence \textsc{high} in the kill)}. Confidence that it reflects a real phenomenon in any tested
habitat: $<5\%$. The programme's question inverts from \emph{when} reversion ignites to \emph{where}
and \emph{whether} reversion robustly exists, and its bottleneck moves from theory to data
trustworthiness. The surviving load-bearing rule: residual reversion and residual half-life are
\emph{smoother-manufactured}---``did price come back?'' is a contaminated question; the valid
discriminator is equilibrium stability (did $\mu^*$ stay put?), and the return-space variance ratio is
the one core statistic that escapes the trap. No transition-morphology variant reopens without a new
mechanism and new evidence.
\end{verdict}

\section{The Provenance and Construction Crisis}
\label{sec:crisis}

\subsection{Provenance is upstream of every statistic}

The kill exposed an inverted ordering: the programme had pursued \emph{when} before establishing
\emph{where}, on data whose causality was unverified. A provenance audit, with disposition rules fixed
before inspection, returned a binding fact: every deployment-domain spread in the data set was a
pre-computed single price column with its constituent legs discarded. With the legs gone, the
hedge-ratio's causality is unrecoverable, and a full-sample hedge ratio---which renders a spread
stationary by hindsight construction---is mechanically indistinguishable from genuine mean reversion.
Consequently \emph{none} of the eight pre-existing spreads survived to the analysis whitelist; all were
contaminated by default. Contamination is the default; trust is earned by a constructive causal proof.

The resolution was a deep cohort of forty-six trusted raw legs across energy, metals, agriculture,
foreign exchange and equities, from which seven spreads are constructible under a causal protocol. For
the first time the programme held instruments whose causality could be proven by construction rather
than assumed.

\subsection{The construction law}

The audit was generalised into a construction-legality standard governing data only, upstream of every
reversion statistic. Its governing question, applied to every derived value at every bar, is whether
that value could have existed using only information available at time $t$. The standard fixes a
hedge-ratio ontology ordered from safest to most assumption-laden---unit ratio (definitional,
zero look-ahead) $\succ$ structural/economic ratio $\succ$ causal rolling regression $\succ$ causal
rolling cointegration $\succ$ dynamic state-space ratio (deferred)---and forbids a full-sample ratio of
any kind as evidence. It requires trailing-window estimation lagged to $t-1$, timestamp (not index)
synchronisation, inner-join alignment, and a level-difference representation for any spread that can
cross zero (logarithmic and ratio statistics are undefined through zero). It treats naive spread
high/low constructed from separately-timed leg extrema as counterfactual states that never existed.

\subsection{A causal estimator that manufactures its own answer}

The first verdict-capable study built seven spreads under this protocol using a causal, lagged
\emph{rolling} ordinary-least-squares hedge ratio on a sixty-bar window, and tested the level-difference
variance ratio against random-walk, GARCH and OU surrogates. Every spread passed a future-injection
causal acceptance test, and the engine was ground-truth validated. The headline result was
nonetheless \emph{inconclusive and construction-defective}. Decompose the differenced spread,
$S_t = A_t - \beta_t B_t$:
\begin{equation}
\Delta S_t \;=\; \underbrace{\big(\Delta A_t - \beta_{t-1}\,\Delta B_t\big)}_{\text{spread innovation}}
\;-\; \underbrace{(\beta_{t-1}-\beta_{t-2})\,B_{t-1}}_{\text{hedge-ratio update noise}\,\times\,\text{trending level}}.
\label{eq:betanoise}
\end{equation}
The second term---the change in the estimated ratio multiplied by the \emph{level} of a trending
leg---was measured at $82$--$97\%$ of $\mathrm{Var}(\Delta S)$ on the failing spreads, with
autocorrelation $\approx 0.81$, injecting a smooth persistent component that drives the variance ratio
far above one (spurious super-diffusion). On synthetic ground truth where the true spread is OU (true
$\VR<1$), the true-ratio variance ratios $[0.97,0.91,0.81,0.65]$ were inflated by the rolling estimator
to $[1.81,3.68,5.46,6.23]$; ablating the update term in \eqref{eq:betanoise} recovered $\VR<1$. The
estimator manufactures the property it measures. The cohort yielded zero admissible confirmations: two
apparent confirmations were a stale-quote artifact ($75.6\%$ flat bars) and a volatility-clustering
near-miss (separates from a random walk at $p=0.005$ but not from a GARCH martingale at $p=0.094$); the
remainder were construction-invalid.

\begin{verdict}{Verdict --- Section \ref{sec:crisis}}
\small
\textbf{No pre-existing spread is admissible} \status{frozen}; the construction law is \status{frozen}.
\textbf{The rolling-OLS-on-levels estimator is killed as inadmissible} for variance-ratio habitat tests
on trending legs (confidence \textsc{high}: synthetic ground truth, ablation, two independent variance
decompositions, a monotone window-length gradient). The load-bearing finding is methodological, not
about the market: a causal time-varying hedge ratio can leak \emph{interpretively}---manufacturing the
measured property---without leaking \emph{temporally}. The deployment-domain reversion premise is thus
\emph{neither confirmed nor validly damaged}; the bottleneck moves from data to construction. The only
artifact-free construction available is the definitional unit ratio, which both fixes the estimator and
supplies a real-data positive control.
\end{verdict}

\section{Apparatus Validation: Natural-Gas Storage Mean Reversion}
\label{sec:apparatus}

\subsection{The positive-control mandate and the unit-ratio restriction}

Every positive control to date had been synthetic, so a negative result was formally indistinguishable
from a too-blunt test. The standing requirement is therefore that, before any further negative is
credible, the apparatus must \emph{confirm} a known, literature- and economically-anchored real edge.
The chosen edge is storage/term-structure mean reversion in a commodity calendar spread---grounded in
the theory of storage \cite{Working1949} and the commodity term-premium literature
\cite{Szymanowska2014}---and the primary instrument is a natural-gas front calendar, $S_t = G_{m1}-G_{m2}$.
Construction is restricted to definitional unit-ratio calendar spreads: the two legs are the same
commodity, so the one-to-one relationship is definitional rather than estimated, the construction
carries zero hedge-ratio degrees of freedom, and it is structurally immune to the artifact of
\eqref{eq:betanoise}.

\subsection{The synthetic calibration gate}

Because the surrogate ensemble was extended with a moving-average ($\mathrm{MA}(1)$) microstructure
null, the calibration was re-earned on synthetic ground truth before any real verdict. All four
controls passed (Table~\ref{tab:gate}). The decisive control is the third: a pure random walk plus
i.i.d.\ level noise (zero mean reversion) produces a variance ratio below one that the random-walk and
GARCH nulls \emph{cannot} absorb---precisely the bid-ask-bounce artifact a calendar spread inherits.
Only the $\mathrm{MA}(1)$-in-increments null discriminates genuine reversion from lag-one
microstructure.

\begin{table}[h]
\centering
\small
\begin{tabular}{@{}p{0.50\linewidth} p{0.24\linewidth} p{0.20\linewidth}@{}}
\toprule
\textbf{Control} & \textbf{Required} & \textbf{Result} \\
\midrule
True seasonal-OU calendar & confirm (power) & confirms \\
True random-walk calendar & null at $\approx 5\%$ & FPR $=3.0\%$ (200 seeds) \\
Random walk $+$ i.i.d.\ noise (zero MR) & $\mathrm{MA}(1)$ must reject & rejected by $\mathrm{MA}(1)$ \\
No-seam invariance (clean OU) & $\VR$ shift $\approx 0$ & shift $=0.000$ \\
\bottomrule
\end{tabular}
\caption{Synthetic calibration gate. The 3.0\% false-positive rate over 200 seeds is below the nominal
5\% and tighter than any single surrogate family because the gate is a conjunction.}
\label{tab:gate}
\end{table}

\subsection{Confirmation}

Table~\ref{tab:confirm} reports the verdict matrix (real minus matched surrogate; multiplicity-corrected
minimum-variance-ratio $p$). The natural-gas calendar beats all three martingale-class nulls at the
fifth-percentile gate, sitting at the surrogate floor. The confirmation survives causal
deseasonalization (the raw confirmation is load-bearing; the deseasonalized strengthening is reported
but not relied upon), is consistent across the open and close sampling points, and---critically---is
corroborated on an \emph{independently constructed} Brent calendar built from raw legs under a
controlled roll rule with no vendor back-adjustment, which confirms at \emph{all} horizons. The
apparatus is also selective: it correctly nulls a refined-products back-month calendar at the book
horizon ($q\leq 20$, $p=0.537$) and recovers its reversion only at the long horizon ($q\geq 60$,
$p=0.005$), exactly the back-month storage profile theory predicts.

\begin{table}[h]
\centering
\footnotesize
\setlength{\tabcolsep}{4.5pt}
\resizebox{\textwidth}{!}{%
\begin{tabular}{@{}l c r l c c c l@{}}
\toprule
\textbf{Instrument} & \textbf{Role} & $n$ & $\VR(2,5,10,20)$ & $p_{\mathrm{RW}}$ & $p_{\mathrm{GARCH}}$ & $p_{\mathrm{MA1}}$ & \textbf{gate} \\
\midrule
NG calendar (daily) & primary & 4969 & $0.92,\,0.83,\,0.61,\,\mathbf{0.45}$ & $0.005$ & $0.025$ & $0.005$ & \textbf{confirm} \\
Brent calendar (built, 60\,min) & control & 22955 & $0.91,\,0.80,\,0.71,\,\mathbf{0.63}$ & $0.005$ & $0.010$ & $0.005$ & \textbf{confirm (all $q$)} \\
RB calendar (daily) & reference & 5170 & $0.996,\,0.993,\,0.983,\,0.979$ & $0.537$ & --- & --- & null (expected) \\
RB long-horizon & --- & 5170 & $q_{40}{=}0.63,\;q_{60}{=}0.49,\;q_{120}{=}0.42$ & $0.005$ & --- & --- & reverts $q{\geq}60$ \\
\bottomrule
\end{tabular}%
}
\caption{Verdict matrix. The natural-gas confirmation is carried at the long horizon ($q=10,20$); the
self-built Brent calendar confirms at all horizons, the opposite of the long-horizon-only compression
signature that vendor back-adjustment would produce.}
\label{tab:confirm}
\end{table}

\subsection{The residual open channel}

An adversarial review returned \emph{holds-with-caveats}. The one artifact the battery does not fully
discriminate is vendor back-adjustment: additive back-adjustment at the monthly roll mechanically
compresses long-horizon variance, producing a declining variance ratio that is indistinguishable from
storage mean reversion under a level-difference statistic. This channel is \emph{partially mitigated}
by the construction-controlled Brent corroboration but \emph{not eliminated} for the natural-gas daily
series specifically.

\begin{verdict}{Verdict --- Section \ref{sec:apparatus}}
\small
\status{conditional survival}: the real-data positive control is \textbf{passed in the conditional
sense}. \textbf{Research win (apparatus works) --- \textsc{medium--high}}: demonstrated power,
selectivity, $3.0\%$ false-positive rate, noise-discrimination, and construction-controlled
corroboration; the synthetic-only prior state is materially improved. \textbf{Empirical win (signal
real) --- \textsc{medium}}: the confirmation is surrogate-relative and robust, but \emph{clean storage
mean reversion} versus \emph{partial vendor-construction artifact} is mitigated, not eliminated.
\textbf{Deployment win --- not claimed}: a positive control is not a deployment target.
\end{verdict}

\section{Rolling-Local Deployability: Persistent but Uneconomic}
\label{sec:rolling}

\subsection{The trader-relevant question and an honest instrument}

Rolling/local estimation is the principal look-ahead and degrees-of-freedom surface---it manufactured
the super-diffusion of \eqref{eq:betanoise}---so the rolling-local analysis was pre-registered under a
binding guard: windows fixed in advance, surrogate-relative under identical extraction, multiplicity
controlled, and local survivability defined as persistence across \emph{multiple disjoint} windows
rather than appearance in one. A first-draft per-window binary flag was falsified pre-registration on
synthetic ground truth (underpowered on genuine OU, firing \emph{higher} on bid-ask bounce, blind to
back-adjustment) and replaced---before any real rolling estimate was computed---by a pooled continuous
statistic. For each window $w$,
\begin{equation}
z_w \;=\; \frac{\VR(20)_w - \overline{\VR(20)}^{\,\mathrm{RW}}_w}{\widehat{\sigma}\big(\VR(20)^{\mathrm{RW}}_w\big)},
\qquad
\text{pooled mean-}z \;=\; \frac{1}{W}\sum_{w=1}^{W} z_w,
\label{eq:poolz}
\end{equation}
with the random-walk band computed in-window so the small-sample variance-ratio bias self-cancels.
Synthetic anchors were frozen first: a random-walk null at $+0.02$ (5--95\% band $[-0.32,+0.41]$),
genuine moderate reversion (OU, $\phi=0.95$) at $-0.33$, and a back-adjustment splice (fraction $0.5$)
at $-0.99$.

\subsection{Results}

The natural-gas signal is statistically persistent (Table~\ref{tab:rolling}): pooled mean-$z=-0.627$
($t=-4.37$, $p\approx 2\times 10^{-4}$) over nineteen disjoint years, with fourteen of nineteen years
below the random-walk median, present both before and after 2020. The global half-life is $12.9$ bars
($\approx 2.6$ weeks, inside the tradeable band). And the structure is reversion-favoring: the signal
\emph{switches off in five years}---all storage-glut / low-scarcity regimes in which the physical
restoring force is weakest, not acute shocks. A uniform back-adjustment artifact would not selectively
switch off in glut years, which is the strongest evidence the persistence is genuine storage mean
reversion.

\begin{table}[h]
\centering
\small
\resizebox{\textwidth}{!}{%
\begin{tabular}{@{}l r r l@{}}
\toprule
\textbf{Read} & \textbf{mean-}$z$ & $n$ & \textbf{note} \\
\midrule
NG yearly (primary) & $\mathbf{-0.627}$ & 19 & $t=-4.37$, $p\approx 2\times10^{-4}$; 14/19 below RW median \\
NG post-2020 & $-0.438$ & 6 & below $-0.32$: recency holds, attenuated \\
Brent (built, no back-adj.) & $-3.54$ & 8 & hourly/off-scale: apparatus control, not NG-level evidence \\
splice frac.\ $0.25$ / $0.5$ & $-0.657$ / $-1.053$ & --- & NG $(-0.627)\approx$ the quarter-strength splice anchor \\
\bottomrule
\end{tabular}%
}
\caption{Rolling-local persistence. The natural-gas pooled distance is statistically indistinguishable
from the quarter-strength splice anchor ($p=0.84$), so genuine reversion versus partial back-adjustment
remains \textsc{medium} confidence.}
\label{tab:rolling}
\end{table}

The economics, however, are decisive. A frozen causal mean-reversion rule (enter at $|z|\geq 1$, exit
at $z=0$ or a twenty-bar stop, round-trip cost $0.003$ spread units) yields, over 200 trades, an
average \emph{gross} capture of $+0.0004$ and an average \emph{net} of $-0.0026$ per trade
(hit-rate $0.53$). The signal is break-even before costs and negative after; clearing the cost floor
would require a gross improvement of roughly seven-fold that no admissible rule-optimization can supply
on this instrument.

\begin{verdict}{Verdict --- Section \ref{sec:rolling}}
\small
\status{persistent-but-uneconomic} \textbf{(confidence \textsc{medium})}. \textbf{Empirical win ---
strengthened, qualified}: the natural-gas calendar sub-diffusion is persistent versus a random walk
(\textsc{high}: $t=-4.37$) and reversion-favoringly structured (switches off in gluts, not shocks);
that it is genuine storage mean reversion versus partial back-adjustment remains \textsc{medium}
($p=0.84$ against the quarter-strength splice). \textbf{Deployment win --- no} (\textsc{high}: the naive
single-instrument book is uneconomic). The deployable object is accordingly redefined from an
\emph{instrument} to a \emph{cost-aware book}; whether a portfolio of individually sub-cost,
leg-netted spread edges can clear costs is the decisive untested question.
\end{verdict}

\section{Corrected Framework and Forward Path}
\label{sec:forward}

\subsection{A rolling-local ontology under a binding guard}

The arc above motivates a change of ontology that is directional, not methodological. The market is
treated as piecewise and regime-local rather than as one global object; the operative question is
\emph{where} and \emph{when} mean reversion exists, not whether it exists universally. Rolling/local
estimation becomes a first-class default but only under the binding guard of
Section~\ref{sec:rolling}---pre-registered windows, surrogate-relative extraction, multiplicity
control, and survivability across multiple disjoint windows---because rolling estimation is the single
most dangerous degrees-of-freedom surface in the programme. Global claims now require cross-habitat
replication plus a mechanism; local conditional claims are the default currency.

\subsection{Deployability as a portfolio property, evaluated after costs}

The deployable object is a \emph{book}, not an instrument: the unit of success is positive
cost-adjusted portfolio expectancy, evaluated after transaction cost, roll, borrow, capacity and
sizing. A spread that mean-reverts but is uneconomic after costs is a non-finding---precisely the
status of the natural-gas calendar. Transaction costs, capacity and risk are first-class research
outputs, not afterthoughts. This prioritization decides \emph{which} questions earn effort; it never
relaxes \emph{how rigorously} a survivor is judged.

\subsection{Open problems}

\begin{description}[leftmargin=1.4em,itemsep=2pt]
  \item[No cross-habitat daily replication.] \status{open} The natural-gas confirmation is a single
  daily instrument; the Brent corroboration is at an off-scale hourly resolution. The unit of evidence
  is survival across independent habitats, not appearance in one.
  \item[Back-adjustment not closed.] \status{open} The natural-gas persistence is statistically
  indistinguishable from a quarter-strength splice; genuine-reversion confidence stays \textsc{medium}
  until the seam is splice-tested or the series is rebuilt from raw legs.
  \item[No controlled-ratio pair.] \status{deferred} The deployment construction proper---pairs and
  cross-asset relative value with a regularized hedge ratio---is untested; the rolling estimator was
  shown to mishandle it, and a regularized ratio must be validated by equilibrium stability rather than
  residual reversion.
  \item[Trend-death detector.] \status{deferred} A trend-death$\,\rightarrow\,$reversion operationalization
  returned no content (a permissive $18$--$20\%$ firing rate, zero of four instruments meeting the
  pre-committed criteria, and anti-thesis continuation in structural growth equities); the thesis is not
  falsified but is untested in the right instrument domain, and any redesign must target a firing rate
  below $5\%$.
  \item[Portfolio economics.] \status{open} The decisive door: whether aggregating sub-cost per-trade
  calendar edges with leg-netting clears costs is the only demonstrated path from ``merely true'' to
  ``deployable''.
\end{description}

\subsection{Three forward hypotheses}

The following are forward directions, not findings; each is open or deferred, carries a reopen trigger,
and inherits the multiplicity and surrogate-relative discipline above.

\begin{description}[leftmargin=1.4em,itemsep=2pt]
  \item[H-Residual-Growth \normalfont(\status{deferred}).] That the expansion/compression dynamics of
  the equilibrium residual around \emph{realized} reversions carry market-specific, conditional
  structure that a matched surrogate does not already produce. This is the legitimate \emph{observational}
  successor to the killed morphology---strictly descriptive, never a detector; any drift toward predicting
  reversion from residual shape is a resurrection and is forbidden without a new, independent
  pre-registration.
  \item[H-Confluence \normalfont(\status{open}).] That a deployable selectivity edge arises not from any
  single signal but from the confluence of multiple independent, surrogate-relative conditions---each
  individually weak, jointly informative for \emph{where not to trade}. Its danger is multiple-comparison
  degrees of freedom; it is admissible only with pre-registered cohorts and cross-habitat
  out-of-sample replication.
  \item[H-Econ-Confluence \normalfont(\status{deferred}, data-gated, long-term only).] The
  economically-conditioned form: that causal, non-price fundamental variables (inventory surprise,
  positioning, signed order flow) condition the confluence---the etiology-conditioned reversion principle
  (inventory/flow $\rightarrow$ reverts; information $\rightarrow$ continues). It is the best-grounded
  mechanism in the programme and is blocked on data acquisition; a price-level conditioning attempt was
  already killed as a selection-on-regime artifact, so the admissible form is surprise- or flow-based.
\end{description}

\subsection{The long-term object}

The terminal ambition is a multi-instrument mean-reversion portfolio engine: a book of independently
cost-clearing, surrogate-validated, causally constructed spread edges across deployment habitats, sized
and risk-managed so that net-of-cost expectancy is positive over a forward hold-out. The path is now
empirically posable for the first time, because the apparatus is validated and one real edge is
confirmed; it is also honestly long, requiring the back-adjustment channel to be closed, replication
across daily habitats, validation of the controlled-ratio construction, and a demonstration that a book
of individually sub-cost edges can net to a cost-clearing aggregate. None of that is done. The posture
is to build that book one validated, cost-aware sleeve at a time---and to stop, and call it a success,
the moment the evidence says the book does not pay.

\section{Conclusion}

Over the interval reported, the programme killed its flagship transition object cleanly and at scale,
diagnosed the structural error that produced it (asking \emph{when} before establishing \emph{where},
on untrustworthy data), rebuilt its data foundation, discovered and removed a construction artifact in
which a causal estimator manufactured its own signal, and---on a corrected, definitional
construction---confirmed for the first time a known real mean-reversion edge with a calibrated,
selective, noise-discriminating apparatus. That same edge, read honestly against transaction costs, is
persistent but uneconomic as a single instrument. The programme therefore holds a research win and an
empirical win and has not earned a deployment win. The remaining path runs through portfolio economics
and causal conditioning variables, never through a resurrected timing object. Throughout, stopping
where the evidence stops is treated as a successful outcome.

\appendix
\section{Definitions and Core Formulae}

\textbf{Level-difference variance ratio.} For a level series $S_t$ that may cross zero,
\[
\VR(q) \;=\; \frac{\mathrm{Var}(S_t - S_{t-q})}{q\,\mathrm{Var}(S_t - S_{t-1})},
\]
with $\VR<1$ indicating sub-diffusion (the mean-reverting direction), $\VR\approx 1$ a martingale, and
$\VR>1$ super-diffusion. Read surrogate-relative (real minus matched null under identical extraction),
never against a fixed constant.

\textbf{Surrogate ensemble.} Random walk, GARCH$(1,1)$, $\mathrm{MA}(1)$ microstructure noise, and OU,
each at identical length and extraction; the gate is the multiplicity-corrected minimum-variance-ratio
$p$ at the fifth percentile.

\textbf{Half-life.} From a causal first-order autoregressive coefficient $\phi$,
$\half = \ln(0.5)/\ln(\phi)$; the tradeable band is $\approx 3$--$40$ bars.

\textbf{Cost hurdle.} A rule is deployable only if the expected net per trade,
$\mathbb{E}[\text{gross capture} - \text{round-trip cost}]$, is positive. For the natural-gas calendar,
gross $+0.0004$ against cost $0.003$ yields net $-0.0026$: uneconomic.

\textbf{Selection-on-deviation.} Conditioning on $|z|\geq\theta$ samples directionally-efficient paths
by construction; it is the complete and sufficient null for the killed pre-reversion morphology.

\begin{thebibliography}{9}
\small
\bibitem{Working1949} Working, H. (1949). The theory of price of storage. \emph{American Economic
Review}, 39(6), 1254--1262.
\bibitem{Szymanowska2014} Szymanowska, M., de Roon, F., Nijman, T., \& van den Goorbergh, R. (2014).
An anatomy of commodity futures risk premia. \emph{Journal of Finance}, 69(1), 453--482.
\bibitem{LoMacKinlay1988} Lo, A.\,W., \& MacKinlay, A.\,C. (1988). Stock market prices do not follow
random walks: evidence from a simple specification test. \emph{Review of Financial Studies}, 1(1),
41--66.
\bibitem{LoMacKinlay1990} Lo, A.\,W., \& MacKinlay, A.\,C. (1990). When are contrarian profits due to
stock market overreaction? \emph{Review of Financial Studies}, 3(2), 175--205.
\bibitem{GeorgeHwang2004} George, T.\,J., \& Hwang, C.-Y. (2004). The 52-week high and momentum
investing. \emph{Journal of Finance}, 59(5), 2145--2176.
\bibitem{JT1993} Jegadeesh, N., \& Titman, S. (1993). Returns to buying winners and selling losers:
implications for stock market efficiency. \emph{Journal of Finance}, 48(1), 65--91.
\bibitem{EngleGranger1987} Engle, R.\,F., \& Granger, C.\,W.\,J. (1987). Co-integration and error
correction: representation, estimation, and testing. \emph{Econometrica}, 55(2), 251--276.
\bibitem{McLeanPontiff2016} McLean, R.\,D., \& Pontiff, J. (2016). Does academic research destroy stock
return predictability? \emph{Journal of Finance}, 71(1), 5--32.
\end{thebibliography}

\end{document}
